\section{Introduction}

Sodium-ion batteries (SIBs) have emerged as a compelling alternative to lithium-ion technology, driven by the rapidly expanding demands of electric vehicles and grid-scale energy storage [1,2]. The natural abundance and low cost of sodium, combined with the compatibility of SIB chemistries with established lithium-ion manufacturing infrastructure, position SIBs for scalable deployment. However, realizing this potential hinges on the development of suitable anode materials that simultaneously deliver high reversible capacity, robust rate capability, and long-term cycling stability — a combination that remains the central bottleneck in SIB advancement [3].

Among anode candidates, carbon-based materials, particularly hard carbon, represent the most mature option but deliver modest specific capacities on the order of 250--300 mAh g$^{-1}$ and exhibit sluggish kinetics at elevated current densities [4]. Alloy-type anodes offer substantially higher theoretical capacities through alloying reactions with sodium. Sn boasts the highest gravimetric capacity at 847 mAh g$^{-1}$ with the lowest working voltage, yet suffers from extreme volume expansion exceeding 400\% that induces pulverization, unstable solid-electrolyte interphase (SEI) formation, low initial Coulombic efficiency, and rapid capacity fading [5]. Sb provides a moderate capacity of 660 mAh g$^{-1}$ with a layered structure that facilitates initial Na$^{+}$ diffusion and relatively favorable rate performance, though it still undergoes severe volume expansion of approximately 290\% and raises toxicity and price volatility concerns [6]. Within this landscape, Bi distinguishes itself through a combination of attributes: a moderate-high gravimetric capacity of 385 mAh g$^{-1}$ paired with an exceptionally high volumetric capacity of approximately 3800 mAh cm$^{-3}$, an intermediate operating potential of around 0.5 V versus Na$^{+}$/Na that mitigates the risk of sodium plating, and large interlayer spacing along the (012) crystallographic plane that promotes Na$^{+}$ diffusion [7,8]. Nevertheless, the practical application of Bi anodes is constrained by its semi-metallic nature, which imparts poor electronic conductivity, and a volumetric expansion of roughly 250\% during the Bi$\leftrightarrow$Na$_{3}$Bi phase transformation [9]. From the perspective of operational safety and volumetric energy density, Bi emerges as the most promising alloy-type anode candidate, provided that its conductivity and volume expansion challenges can be simultaneously addressed.

Carbon coating strategies have been widely employed to enhance conductivity and buffer volume changes in alloy-type anodes, but the physical adhesion between carbon and active material is often limited and prone to delamination over extended cycling [10]. MXenes, and Ti$_{3}$C$_{2}$T$_{x}$ in particular, offer a uniquely advantageous scaffold for Bi-based composites: metallic conductivity approaching 10,000 S cm$^{-1}$ enables rapid electron transport, hydrophilic surface terminations (---O, ---OH, ---F) provide strong interfacial binding to active species, and the mechanically robust layered structure is inherently capable of accommodating substantial volume fluctuations during sodiation and desodiation [11,12]. Recent studies have explored the combination of Bi with MXene nanosheets for SIB anodes [13], and a MXene@Bi$_{2}$S$_{3}$/Mo$_{7}$S$_{8}$ sandwich heterostructure was shown to deliver 474.9 mAh g$^{-1}$ at 5.0 A g$^{-1}$ with stable cycling over 1,400 cycles, underscoring the potential of sandwich-type MXene configurations [14]. Despite these advances, a structurally simple binary Bi/MXene composite that simultaneously achieves high rate capability and ultra-long cycling stability — without relying on sulfidation, carbon coatings, or multi-component heterostructure engineering — has yet to be demonstrated.

In this work, we report a Bi/MXene composite synthesized via a facile hydrothermal approach, in which Bi nanoparticles are confined between Ti$_{3}$C$_{2}$T$_{x}$ MXene layers in a distinctive sandwich architecture. The MXene component serves a dual function: it establishes a continuous metallic conductive network that enables fast electron transport, while mechanically constraining the volumetric expansion of Bi during sodiation to suppress pulverization and preserve electrode integrity. The resulting Bi/MXene composite delivers a specific capacity of 327 mAh g$^{-1}$ at a high current density of 10 A g$^{-1}$, with an initial Coulombic efficiency of 88\%. Notably, the electrode retains 311 mAh g$^{-1}$ after 5,000 cycles at 2 A g$^{-1}$, corresponding to 92\% capacity retention. These metrics collectively demonstrate that the Bi/MXene sandwich architecture can simultaneously achieve high rate capability, high reversible capacity, and long-term durability — effectively breaking the rate-capacity-stability trade-off that has constrained conventional Bi-based anodes. The underlying electrochemical mechanism is elucidated through detailed kinetics analysis combining cyclic voltammetry at varied scan rates, electrochemical impedance spectroscopy, and galvanostatic intermittent titration technique measurements, complemented by post-cycling scanning electron microscopy and ex-situ X-ray diffraction and X-ray photoelectron spectroscopy.
